\documentclass[12pt,a4paper]{report}
\usepackage[]{amsmath,amsthm,amssymb,amscd}
\usepackage[latin1]{inputenc}
%\usepackage{fullpage}

\topmargin 0pt
\advance \topmargin by -\headheight
\advance \topmargin by -\headsep
     
\textheight 246mm
     
\oddsidemargin 0pt
\evensidemargin \oddsidemargin
\marginparwidth 0.5in
     
\textwidth 159mm


%               Theorem Environments
\theoremstyle{definition}
\newtheorem{ex}{}
\newcommand{\pd}[1]{\frac{\partial}{\partial #1}} %\pd{x}
%               Subquestions numbered with letters
\renewcommand{\theenumi}{\textbf{\alph{enumi}}}

%               Maths definitions

\newcommand{\NN}{\ensuremath{\mathbb N}}
\newcommand{\ZZ}{\ensuremath{\mathbb Z}}
\newcommand{\CC}{\ensuremath{\mathbb C}}
\newcommand{\RR}{\ensuremath{\mathbb R}}
\newcommand{\QQ}{\ensuremath{\mathbb Q}}
\newcommand{\TT}{\ensuremath{\mathbb T}}
\newcommand{\g}{\ensuremath{\frak{g}}}
\newcommand{\h}{\ensuremath{\frak{h}}}
\newcommand{\ft}{\ensuremath{\frak{t}}}
\newcommand{\cB}{\mathcal{B}}
\newcommand{\cN}{\mathcal{N}}
\newcommand{\cL}{\mathcal{L}}
\newcommand{\cD}{\mathcal{D}}
%\newcommand{\pd}[1]{\frac{\partial}{\partial #1}} %\pd{x}
%\newcommand{\pd}[1]{\frac{\partial}{\partial #1}} %\pd{x}

\begin{document}
\noindent
{\sffamily\large Marco Zambon 
\hfill   a discutir:   Mi\'ercoles, 07/03/2012\\Geometr\'ia III, curso 2011-12}
 

\noindent
\hrulefill{}\\
\begin{center}\bfseries\large\textsf{Hoja 2    \vspace{-2mm}}
\end{center}

\noindent
%\hrulefill{}\\

\begin{ex} 
Para cada $\alpha \in [0,1]$, considera la aplicaci\'on
$$\Phi \colon \RR\to S^1\times S^1, t\mapsto (e^{2\pi i t}, e^{2\pi i t \alpha}),$$
donde $S^1:=\{z\in \CC: |z|=1\}$. Supongamos que $\alpha  \notin \QQ$.
\begin{enumerate} 
\item Demuestra que   $\Phi$ es inyectiva y continua.
\item Demuestra que $\Phi \colon \RR \to \Phi(\RR)$ no es un homeomorfismo. Por qu\'e no hay ninguna contradici\'on con la pregunta 4a de la hoja 1? (Aqu\'i $\Phi(\RR)$ tiene la topolog\'ia subespacio inducida por el toro $S^1\times S^1$). 
\item  Es $\Phi(\RR)$ es una variedad topol\'ogica o no?
\end{enumerate}
\end{ex}


\begin{ex}
Sea $K$ una superficie topol\'ogica compacta, conexa y con borde.
Demuestra que el borde $\partial K$ tiene un numero finito de componentes conexas.
\end{ex} 
 
\begin{ex}  
Demuestra para el espacio proyectivo $P_2(\RR)$, la banda de M\"obius $M$, la botella de Klein $K$ y el toro $\TT=S^1\times S^1$:
\begin{enumerate}
\item  $\chi(P_2(\RR))=1$
\item  $\chi(K)=0$
\item $\chi(M)=0$
\item $\chi(\TT)=0$
\end{enumerate}
\textbf{Consejo:}  Para a. puedes utilizar el hecho que $\chi(S^2)=2$
% y para b. utiliza la formula $\chi(S_1\sharp S_2)=\chi(S_1)+\chi(S_2)-2$.
\end{ex}
 

\begin{ex}  
A cual superficie en la lista (de superficies compactas, conexas, con borde) es homeomorfa la suma conexa
$$\TT \;\;\sharp\;\; C,$$
donde $\TT=S^1\times S^1$ y $C=S^1\times [0,1]$?
\end{ex} 
 

\begin{ex}[{\bf Por escrito}]
Demuestra que si $M,N,L$ so variedades diferenciales y $f \colon M \to N$, $g \colon N \to L$ aplicaciones diferenciables, entonces $g \circ f \colon M \to L$ es  una aplicaci\'on diferenciable.
\end{ex} 
 



\begin{ex}
Encuentra una estructura de variedad diferencial en $$Q:=\{(x,y):x\in[-1,1], y\in\{-1,1\}\}\cup \{(x,y): x\in\{-1,1\},y\in[-1,1]\}$$
(el borde de un cuadrado) con la topolog\'ia subespacio inducida por $\RR^2$.
\end{ex} 


\begin{ex} Sea $M$ una variedad diferencial conexa y $p,q$ puntos de $M$. Demuestra que  existe una aplicaci\'on diferenciable $c \colon [0,1] \to M$ tal que $c(0)=p,c(1)=q$.\\
{\bf Observaci\'on:} Como $[0,1]$ no es una variedad diferencial segundo la definici\'on dada  en clase (porque tiene borde), hay que definir cuando una aplicaci\'on  $c \colon [0,1] \to M$ es diferenciable: lo es cuando existe un $\epsilon >0$ y una aplicaci\'on diferenciable $\tilde{c} \colon (-\epsilon, 1+\epsilon) \to M$ tal que 
$\tilde{c}|_{[0,1]}=c$.\\
{\bf Consejo:} Puedes utilizar, sin demonstraci\'on, el hecho que una variedad topol\'ogia conexa es automaticamente conexa por caminos.
 \end{ex} 
 



 \end{document}
 %%%%%%%
 \begin{ex}

\end{ex}

 \begin{ex}
\begin{enumerate}
\item 
\end{enumerate}
\end{ex}
 
 
 
 